\documentclass[11pt,a4paper]{ctexart}
\usepackage[margin=2.45cm]{geometry}
\usepackage{amsmath,amssymb,amsthm,mathtools,bm}
\usepackage{xcolor,booktabs,longtable,array,enumitem,hyperref}
\hypersetup{colorlinks=true,linkcolor=blue!55!black,urlcolor=blue!55!black,citecolor=blue!55!black}
\setlist{nosep}

\newtheorem{theorem}{定理}[section]
\newtheorem{lemma}[theorem]{引理}
\newtheorem{proposition}[theorem]{命题}
\theoremstyle{definition}
\newtheorem{definition}[theorem]{定义}
\newtheorem{remark}[theorem]{说明}

\newcommand{\F}{\mathbb F}
\newcommand{\T}{\mathbb T}
\newcommand{\eps}{\varepsilon}
\newcommand{\LL}{\mathcal L}
\newcommand{\E}{\mathcal E}
\newcommand{\K}{\mathcal K}
\newcommand{\OPEN}{\textcolor{red!70!black}{\textsf{OPEN}}}
\newcommand{\REDUCED}{\textcolor{orange!80!black}{\textsf{REDUCED}}}
\newcommand{\PROVED}{\textcolor{green!45!black}{\textsf{PROVED}}}
\newcommand{\EXTERNAL}{\textcolor{blue!70!black}{\textsf{EXTERNAL}}}

\title{GFT--Goldbach 的 G1/G2/G3 归约\\
\large 条件性主定理、版本消歧与可形式化状态审计}
\author{依据八份研究日志重建的审计稿}
\date{2026年8月28日}

\begin{document}
\maketitle

\begin{abstract}
本文把多轮聊天中不断变动的记号冻结为一个可审计的三模块归约。
严格结论不是“Goldbach 已证明”，而是：在一个标准的加权圆法主弧输入下，
若 G1 给出 fixed-$N$、pre-Bessel、保留 minor-arc 核和变量来源的无损编译，
且 G2、G3 分别控制编译所得的 balanced 与 polarized residual，
则所有充分大的偶数都是两个素数之和。
本文证明这一条件蕴含，解释 G1 的确切内容，消除“G3 有三块”与后来
G1/G2/G3 命名之间的歧义，并给出截至日志末端的保守状态：
G1 开放；G2 已缩到 DCCS/KPA8 类型拼接与高矩回代，但仍开放；
G3 的两素数载体与大部分有限格路由已有强后端，中央单素数 outer-prime gate
仍开放。任何聊天中的 \textsf{PROVED} 标签，若没有在本文中给出完整陈述与证明，
均不被提升为无条件定理。
\end{abstract}

\section{结论先行：这里的“Goldbach = G1+G2+G3”是什么意思}

符号“$=$”不是数值等式，也不是三个彼此独立的充分条件；它是一个串行依赖：
\[
\boxed{
\mathrm{G1}\quad\Longrightarrow\quad
\bigl(\mathrm{G2}\oplus\mathrm{G3}\bigr)
\quad\Longrightarrow\quad
\E_{\mathfrak m}(N)=o(N)
\quad\Longrightarrow\quad
\text{Goldbach for sufficiently large even }N.}
\]
G1 负责证明真实 fixed-$N$ minor arc \emph{确实到达}局部核；G2/G3 才负责估计这些核。
因此即使局部估计很强，没有 G1 也不能回到 Goldbach；反之只有 G1 而没有
G2/G3，也只是把困难改写了一次。

\paragraph{审计判决。}
上述\emph{条件蕴含}可以严格证明；三个条件本身在所给八份日志中没有全部证明。
故本文不声称强 Goldbach 猜想已经解决。

\section{从加权表示函数到 Goldbach}

令
\[
\Lambda(n)=\begin{cases}\log p,&n=p^k,\ k\ge1,\\0,&\text{otherwise},\end{cases}
\qquad
\vartheta(n)=\begin{cases}\log p,&n=p,\\0,&\text{otherwise}.
\end{cases}
\]
固定非负 $W\in C_c^\infty((0,1))$，并假设
\(
\mathfrak J_W:=\int_0^1W(t)W(1-t)\,dt>0.
\)
定义
\[
R_{\Lambda,W}(N)
=\sum_{m+n=N}\Lambda(m)\Lambda(n)W(m/N)W(n/N),
\]
并以 $R_{\vartheta,W}(N)$ 类似定义素数--素数表示函数。

\begin{lemma}[真素数与素数幂的差]
对 $N\ge3$，有
\[
R_{\Lambda,W}(N)-R_{\vartheta,W}(N)
=O_W\!\left(N^{1/2}(\log N)^2\right).
\]
\end{lemma}

\begin{proof}
差项中至少有一侧是 proper prime power $p^k$，$k\ge2$。
不超过 $N$ 的这类数有 $O(N^{1/2})$ 个；每个乘积权至多 $(\log N)^2$，
且另一变量由 $m+n=N$ 唯一确定。交换两侧只改变绝对常数。
\end{proof}

对偶数 $N$，记标准 Goldbach 奇异级数
\[
\mathfrak S(N)=2C_2\prod_{\substack{p\mid N\\p>2}}\frac{p-1}{p-2}>0,
\qquad
C_2=\prod_{p>2}\left(1-\frac1{(p-1)^2}\right)>0.
\]

\begin{definition}[主弧输入]
记 $L=\log N$。本文把下式作为经典圆法的外部输入：对任意固定 $A>0$，
\[
\int_{\mathfrak M}S_\Lambda(\alpha)^2e(-N\alpha)\,d\alpha
=\mathfrak S(N)\mathfrak J_WN+O_{A,W}(NL^{-A}),
\tag{MA}
\]
其中 $S_\Lambda$ 含有与 $W$ 一致的光滑权。
\end{definition}

这一定义只冻结本文需要的接口；它不把任何尚未证明的 GFT 编译步骤藏入“标准主弧”。

\section{G1、G2、G3 的稳定定义}

\subsection{G1：fixed-$N$ pre-Bessel PB--MRDET 编译器}

设 $\E_{\rm hard}(N)$ 为做完主弧、Type-I/easy 分支、显式 endpoint 修正后
仍需处理的 minor-arc nonzero-orbit 项。G1 是下列\emph{恒等式级}陈述。

\begin{definition}[G1]
对每个充分大的偶数 $N$ 和任意 $A>0$，存在有限指标集
$\LL_{\rm bal}(N)$、$\LL_{\rm pol}(N)$、系数 $c_\lambda,d_\mu$ 及
type-correct 局部核 $\K_\lambda^{\rm bal}(N)$、
$\K_\mu^{\rm pol}(N)$，使得
\begin{align}
\E_{\rm hard}(N)
={}&\sum_{\lambda\in\LL_{\rm bal}(N)}c_\lambda
\K_\lambda^{\rm bal}(N)
+\sum_{\mu\in\LL_{\rm pol}(N)}d_\mu
\K_\mu^{\rm pol}(N)
+O_A(NL^{-A}),                                           \tag{G1.1}\\
\sum_\lambda|c_\lambda|+\sum_\mu|d_\mu|
\ll{}& L^{C},                                             \tag{G1.2}
\end{align}
其中每个核保留 $N$、minor-arc Fourier 核、dyadic support、
prime/factor origin、gcd 分支、GFT/Mellin charge 与 Poisson prefactor。
\end{definition}

G1 的关键不是得到一个形似
\(Abf+Bdh=N\) 的式子，而是证明 minor-arc 截断没有被非法替换成全圆正交性。
一般地
\[
\int_{\mathfrak m}e\bigl((\Phi-N)\alpha\bigr)\,d\alpha
=\widehat{1_{\mathfrak m}}(N-\Phi),
\]
而不是 $1_{\Phi=N}$。任何合法 G1 都必须运输这个 Fourier 核，或证明一个
具有可控 Wiener 范数的等价 shifted family。

\subsection{G2：balanced weighted four-cycle}

\begin{definition}[G2]
对 G1 实际产生的 balanced kernels 及其实际系数，有一致重构界
\[
\left|\sum_{\lambda\in\LL_{\rm bal}(N)}
c_\lambda\K_\lambda^{\rm bal}(N)\right|
\ll_{A,W}NL^{-A}.                                         \tag{G2}
\]
这里的 G2 是 reconstructed/weighted 陈述；它不要求每个冻结的 GFT mode
分别满足更强的点态界。
\end{definition}

在日志的局部记号中，G2 是 WPMR/PR4 的 weighted finite-field four-cycle，
目标通常等价于某个 centered covariance 的
\(S_4^4\ll p^{3/2+o(1)}\) 型界。该等价必须连同全部 $p,Q=p-1$
因子逐式证明，不能仅由量纲猜测。

\subsection{G3：polarized/unbalanced finite-cell exhaustiveness}

\begin{definition}[G3]
对 G1 实际产生的所有 polarized residual，有
\[
\left|\sum_{\mu\in\LL_{\rm pol}(N)}
d_\mu\K_\mu^{\rm pol}(N)\right|
\ll_{A,W}NL^{-A},                                         \tag{G3}
\]
并且 $\LL_{\rm pol}(N)$ 的 cell manifest 是穷尽、互斥且保留变量来源的。
\end{definition}

日志中 G3 的因子几何被压缩为三类：
\[
\begin{array}{c|c|c}
\text{cell}&\text{尺度条件（示意）}&\text{载体}\\ \hline
\mathrm{P1}&\lambda_1>2/3& p\\
\mathrm{P2a}&\lambda_1\le2/3,\ \lambda_2>1/3&p_1p_2\\
\mathrm{P2b}&\lambda_1\le2/3,\ \lambda_2\le1/3,
\ \lambda_1+\lambda_2>5/6&p_1p_2.
\end{array}
\tag{3.1}
\]
这些阈值只有在 packing theorem 与原始 support 条件被完整写入后才可使用。

\section{条件性 Goldbach 主定理}

\begin{theorem}[G1+G2+G3 推出充分大二元 Goldbach]
假设主弧输入 \textup{(MA)}、所有 easy/Type-I 与 prime-power 修正为 $o(N)$，
并且 G1、G2、G3 对每个充分大的偶数 $N$ 一致成立。则存在 $N_0$，使每个
偶数 $N\ge N_0$ 都可写成两个素数之和。
\end{theorem}

\begin{proof}
由 G1.1 及 G2、G3，
\[
\E_{\rm hard}(N)=O_A(NL^{-A}).
\]
加回已经控制的 easy/Type-I 分支，取 $A$ 足够大，得到完整 minor arcs
$\E_{\mathfrak m}(N)=o(N)$。与 (MA) 相加，
\[
R_{\Lambda,W}(N)=\mathfrak S(N)\mathfrak J_WN+o(N).
\]
对偶数 $N$，$\mathfrak S(N)\mathfrak J_W>0$，故右侧对充分大 $N$ 为正且
为 $\gg_WN$。由素数幂差引理，
\[
R_{\vartheta,W}(N)
=R_{\Lambda,W}(N)+O_W(N^{1/2}L^2)>0
\]
对充分大偶数成立。于是至少存在素数 $p_1,p_2$，满足 $p_1+p_2=N$。
\end{proof}

\begin{remark}[“充分大”与完整猜想]
上一定理本身只给 sufficiently-large Goldbach。要得到“每个 $N\ge4$”，
还需对有限区间 $4\le N<N_0$ 作独立验证；八份日志没有提供一个可计算的
$N_0$ 与相应验证证书，因此本文不把两者合并为完整无条件证明。
\end{remark}

\section{命名漂移：所谓“三块 G3”应怎样理解}

日志中出现过两套分类，若不分开会产生表面矛盾。

\begin{center}
\begin{tabular}{p{0.27\textwidth}p{0.30\textwidth}p{0.30\textwidth}}
\toprule
分类层级 & 三分法 & 后来的 G1/G2/G3 命名\\
\midrule
Local SOMK 轨道
& zero orbit；balanced nonzero；polarized nonzero
& zero orbit 为前置 bookkeeping；balanced 成为 G2；polarized 成为 G3\\
全局剩余障碍
& 不适用
& G1=PB--MRDET；G2=balanced WPMR/PR4；G3=polarized finite cells\\
G3 内部 finite cells
& P1；P2a；P2b
& 仍统称 G3，不再与 zero/balanced 混合\\
\bottomrule
\end{tabular}
\end{center}

因此“local SOMK 的三块已有两块完成、最后一块只余渣滓”是较早口径；
后来把最后的 polarized 块单独命名为 G3。本文采用后一口径。

\section{截至日志末端的严格状态}

\begin{longtable}{p{0.25\textwidth}p{0.14\textwidth}p{0.52\textwidth}}
\toprule
模块 & 状态 & 审计说明\\
\midrule
\endfirsthead
\toprule
模块 & 状态 & 审计说明\\
\midrule
\endhead
标准加权主弧与奇异级数正性 & \EXTERNAL
& 可作为经典外部输入；不包含 GFT 的 minor-arc 编译。\\
有限循环群 DFT/Mellin 正交、sum--product 有限纤维
& \PROVED & 可在有限和层面完全形式化；见配套 toolkit 稿。\\
G1 / PB--MRDET
& \OPEN & 日志给出 HWAC/shifted-MRDET 等候选链，但没有一个把真实
minor-arc cutoff、全部叶子、误差与类型同时闭合的最终 theorem。\\
G2 / balanced WPMR--PR4
& \REDUCED & 旧 affine replay 因 Mellin 共轭方向错误而失效；新路线到
CRAF/KPA8，最后停在 DCCS exact operator factorization 与 $m$-$L^8$
资源的同一 normalization 回代。不得标为 closed。\\
G3 / 2P cells
& \REDUCED & 固定 semiprime modulus 有 Pascadi 的强后端；但从日志核到
外层加权总和仍需逐 cell 检查 coprimality、长度、模分解与 coefficient norms。\\
G3 / PCRS typing
& \REDUCED & “Poisson raw legs 自动来自同一 $X^{1/3}$ cofactor”的强版本
已被日志自己否定；较弱的 carrier isolation 是可行的 origin invariant，但
仍需嵌入完整 G1 manifest。\\
G3 / central 1P outer-prime
& \OPEN & 被缩到 BP trace pushforward 的 weighted square-core paucity；
现有 polynomial-square paucity 是未加权 $P(n)$ 计数，不能直接控制 GFT 权 $a_R$。\\
Binary Goldbach
& \OPEN & 条件归约成立，但三个全局 obligation 未全部证明。\\
\bottomrule
\end{longtable}

\section{G1 的精确定义：它到底要证明什么}

G1 可以用以下 proof-object 表达。每个 leaf $\lambda$ 至少包含：
\[
\lambda=(N,\,\text{arc kernel},\,\text{dyadic boxes},\,
\text{origin labels},\,\text{gcd branch},\,\text{GFT charges},\,
\text{Poisson scales}).
\]
一个 Lean-friendly 的 G1 不应只返回一个复数等式，而应返回：
\begin{enumerate}[label=\textup{(G1-\arabic*)}]
\item 一个有限 leaf 列表及其互斥/穷尽证明；
\item 从原 minor-arc 和到 leaf 和的精确恒等式或有显式误差的等式；
\item 每次变量替换的双射、定义域与非零分母证明；
\item 每个 Fourier/Poisson/Mellin 变换的归一化常数；
\item 绝不把 $\widehat{1_{\mathfrak m}}(N-\Phi)$ 偷换成 $1_{\Phi=N}$；
\item zero orbit 被删除还是并入奇异级数的互斥 bookkeeping；
\item 输出 kernel 满足 G2 或 G3 theorem 的全部 hypotheses 的证明。
\end{enumerate}
这就是 G1 比“代数上写出一个 determinant equation”强得多的地方。

\section{Lean 形式化边界}

“Lean-friendly”在本文中意为：每个有限域恒等式可被拆成有限类型、有限和、
显式分母 guard 与范数不等式；并不声称本文已经附带可编译 Lean 源码。
推荐的形式化层次为：
\begin{enumerate}
\item \textbf{Finite algebra layer}：$\F_p$、字符正交、DFT/Mellin、有限矩阵与
Schatten 范数；可完全形式化。
\item \textbf{Exact compiler layer}：Vaughan identities、dyadic partitions、
origin-labelled AST 与每个 leaf 的语义保持；这是 G1 的核心。
\item \textbf{Analytic import layer}：把 Blomer--Pascadi、Pascadi 等论文定理
封装成带全部假设的结构，禁止“只导入结论指数”。
\item \textbf{Asymptotic layer}：所有 $p^{o(1)}$ 改写成
$\forall\eps>0,\exists C_\eps$ 的量词陈述。
\end{enumerate}

\section{外部定理的适用范围警告}

\begin{itemize}
\item Blomer--Pascadi 的平方根临界 saving 是固定模数、双线性 Kloosterman 形式的
定理；它不自动提供 outer-prime 平均。
\item Pascadi 的 non-abelian amplification 在有限群不可约表示上成立，且分母
来自字符在正规子群上的质量；“multiplicity-one”本身不推出任意 coefficient
operator diagonal。
\item Kable 的 Legendre/Soto--Andrade 正交性位于带端点质量的
$\ell^2(\F_q,m)$；迁移到本文的 counting measure 必须显式处理权重。
\item de la Bret\`eche--Wang--Xu 的 quadratic polynomial square-paucity
不包含本文的 GFT pushforward 权 $a_R$。加权版本是新的 obligation。
\end{itemize}

\section{最终依赖图与下一步}

\[
\boxed{
\begin{array}{c}
\text{standard major arcs + exact endpoint algebra}\\
\Downarrow\\
\mathrm{G1:\ fixed\text{-}N\ lossless\ PB\!\!-
MRDET}\quad(\OPEN)\\
\Downarrow\\
\begin{matrix}
\mathrm{G2:\ balanced\ WPMR/PR4}\ (\REDUCED)
&+&
\mathrm{G3:\ polarized\ cells}\ (\REDUCED/\OPEN)
\end{matrix}\\
\Downarrow\\
\E_{\mathfrak m}(N)=o(N)\\
\Downarrow\\
R_{\vartheta,W}(N)>0.
\end{array}}
\]

短期最小任务不是再宣布一个新的总定理，而是：先完成 G2 的 DCCS
kernel-by-kernel identity 与 normalization replay；并行的 G3 最小解析目标是
central 1P weighted square-core BPTPE。G1 仍是最终把局部成果接回 fixed-$N$
Goldbach 的最大接口。

\begin{thebibliography}{9}
\bibitem{Kable2002}
A.~C. Kable,
\emph{Legendre sums, Soto--Andrade sums and Kloosterman sums},
Pacific J. Math. 206 (2002), 139--157.
\url{https://msp.org/pjm/2002/206-1/pjm-v206-n1-p09-s.pdf}

\bibitem{Pascadi2026}
A. Pascadi,
\emph{Non-abelian amplification and bilinear forms with Kloosterman sums},
arXiv:2511.08445v2 (2026).
\url{https://arxiv.org/abs/2511.08445}

\bibitem{BP2026}
V. Blomer and A. Pascadi,
\emph{Bilinear forms with Kloosterman sums via quadratic characters},
arXiv:2607.24311v1 (2026).
\url{https://arxiv.org/abs/2607.24311}

\bibitem{dBWX2026}
R. de la Bret\`eche, V.~Y. Wang and M.~W. Xu,
\emph{Random Multiplicative Functions and Making Squares from Polynomial Values},
arXiv:2607.06398v1 (2026).
\url{https://arxiv.org/abs/2607.06398}
\end{thebibliography}

\end{document}
